\begin{thebibliography}{100}
\bibitem{ifrah2000universal}
Georges Ifrah, \textit{The Universal History of Numbers: From Prehistory to the Invention of the Computer}, Wiley, 2000.

\bibitem{ore1990number}
Oystein Ore, \textit{Number Theory and Its History}, Dover Publications, 1990.

\bibitem{kline1990mathematics}
Morris Kline, \textit{Mathematics: The Loss of Certainty}, Oxford University Press, 1990.

\bibitem{ross2003elementary}
Kenneth A. Ross, Charles R. B. Wright, \textit{Elementary Number Theory}, 6th Edition, Pearson, 2003.

\bibitem{strang1993introduction}
Gilbert Strang, \textit{Introduction to Linear Algebra}, Wellesley-Cambridge Press, 1993.

\bibitem{knuth1997art}
Donald E. Knuth, \textit{The Art of Computer Programming, Volume 2: Seminumerical Algorithms}, Addison-Wesley, 1997.



\bibitem[Amdahl '64]{Amdahl} G. M. Amdahl, G. A. Blaauw, F. P. Brooks, Jr. (1964)  Architecture of the IBM System/360, {\em IBM Journal of Research and Development}, {\bf 8}(2):87--101.

\bibitem[Alfke '05]{Alfke} Peter Alfke: ``Metastable Recovery in Virtex-II Pro FPGAs", {\em Application Note: Virtex-II Pro Family}, {\tt http://www.xilinx.com/support/documentation/application\_notes/xapp094.pdf}, XILINX, 2005.

\bibitem[Andonie '95]{} R\v azvan Andonie, Ilie G\^{a}rbacea: {\em Algoritmi fundamentali. O perspectiv\v a $C^{++}$}, Ed. Libris, Cluj-Napoca, 1995. (in Roumanian)

\bibitem[Ajtai '83]{Ajtai} M. Ajtai, et al.: ``An O(n log n) sorting network", Proc. 15th Ann. ACM Symp. on Theory of Computing, Boston, Mass., 1983

\bibitem[Asanovic '06]{Asanovic06} K. Asanovic, {\em et al.} (2006) {\em The landscape of parallel computing research: A view from Berkeley}, Technical Report EECS-2006-183, Berkeley University, Berkeley, CA. At: \url{https://www.eecs.berkeley.edu/Pubs/TechRpts/2006/EECS-2006-183.pdf}

\bibitem[Asanovic '09]{Asanovic09} K. Asanovic,  {\em et al.} (2009) {\em A View of the Parallel Computing Landscape}, Communications of the
    ACM, {\bf 52}(10):56--67.

\bibitem[Backus '78]{Backus} J. Backus (1978) Can programming be liberated from the von neumann style? a functional style and its algebra of programs, {\em Communications of the ACM}, {\bf 21}(8):613--641.

\bibitem[Batcher '68]{Batcher} K. E. Batcher: ``Sorting networks and their applications", in {\em Proc. AFIPS Spring Joint Computer Conference}, vol. 32, 1968.

\bibitem[Benes '68]{Benes} V\' aclav E. Bene\v s: {\em Mathematical Theory of Connecting Networks and Telephone Traffic}. New York: Academic, 1968.

%\bibitem[Blakeslee '79]{} T. R. Blakeslee: {\em Digital Design with Standard MSI and LSI}, John Wiley \& Sons, 1979.

%\bibitem[Booth '67]{} T. L. Booth: {\em Sequential Machines and Automata Theory}, John Wiley \& Sons, Inc., 1967.

%\bibitem[Bremermann '62]{} H. G. Bremermann: ``Optimization through Evolution and Recombination", in {\em Self-Organizing Systems}, ed.: M. C. Yovits, S. Cameron, Washington DC, Spartan, 1962.

%\bibitem[Calude '82]{} Cristian Calude: {\em Complexitatea calculului. Aspecte calitative} (The Complexity of Computation. Qualitative Aspects), Ed. Stiintifica si Enciclopedica, Bucuresti, 1982.

%\bibitem[Calude '94]{} Cristian Calude: {\em Information and Randomness}, Springer-Verlag, 1994.

\bibitem[Casti '92]{Casti} John L. Casti: {\em Reality Rules: II. Picturing the World in Mathematics - The Frontier}, John Wiley \& Sons, Inc., 1992.

%\bibitem[Cavanagh '07]{Cavanagh} Joseph Cavanagh: {\em Sequential Logic. Analysis and Synthesis}, CRC Taylor \& Francis, 2007.

\bibitem[Chaitin '66]{Chaitin66} Gregory Chaitin: ``On the Length of Programs for Computing Binary Sequences", {\em J. of the ACM}, Oct., 1966.
\bibitem[Chaitin '70]{Chaitin70} Gregory Chaitin: ``On the Difficulty of Computation", in {\em IEEE Transactions of Information Theory}, ian. 1970.
\bibitem[Chaitin '77]{Chaitin '77} Gregory Chaitin: ``Algorithmic Information Theory",  in {\em IBM J. Res. Develop.}, Iulie, 1977.
%\bibitem[Chaitin '87]{} Gregory Chaitin: {\em Algorithmic Information Theory}, Cambridge University Press, 1987.
%\bibitem[Chaitin '90]{} Gregory Chaitin: {\em Information, Randomness and Incompletness}, World  Scientific,1990.
%\bibitem[Chaitin '94]{} Gregory Chaitin: {\em The Limits of Mathematics IV}, IBM Research Report RC 19671, e-print chaodyn/9407009, July 1994.
%\bibitem[Chaitin '06]{chaitin06} Gregory Chaitin: ``The Limit of Rason", in {\em Scientific American}, Martie, 2006.


\bibitem[Chomsky '56]{Chomsky} Noam Chomsky, ``Three Models for the Description of Languages", {\em IEEE
Trans. on Information Theory}, 2:3 , 1956.
\bibitem[Chomsky '59]{Chomsky1} Noam Chomsky, ``On Certain Formal Properties of Grammars", {\em
Information and Control}, 2:2, 1959.
\bibitem[Chomsky '63]{Chomsky2} Noam Chomsky, ``Formal Properties of Grammars", {\em Handbook of
Mathematical Psychology}, Wiley, New-York, 1963.

\bibitem[Church '36]{Church} Alonzo Church: ``An Unsolvable Problem of Elementary
Number Theory", in {\em American Journal of Mathematics}, vol. 58,
pag. 345-363,
 1936.

%\bibitem[Clare '72]{} C. Clare: {\em Designing Logic Systems Using State Machines}, Mc Graw-Hill, Inc., 1972.

\bibitem[Cormen '90]{Cormen} Thomas H. Cormen, Charles E. Leiserson, Donsld R. Rivest: {\em Introduction to Algorithms}, MIT Press, 1990.

%\bibitem[Dasc\v alu '98]{Dascalu} Monica Dasc\v alu, Eduard Fran\c ti, Gheorghe \c Stefan: ``Modeling Production with Artificial Societies: the Emergence of Social Structure", in S. Bandini, R. Serra, F. Suggi Liverani (Eds.): {\em Cellular Automata: Research Towars Industry. ACRI '98 - Proceedings of the Third Conference on Cellular Automata for Research and Industry}, Trieste, 7 - 9 October 1998, Springer Verlag, 1998. p 218 - 229.

%\bibitem[Dasc\v alu '98a]{Dascalua} Monica Dasc\v alu, Eduard Fran\c ti, Gheorghe \c Stefan: ``Artificial Societies: a New Paradigm for Complex Systems' Modeling", in {\em IFAC Conference on Supplemental Ways for Improving International Stability - SWIIIS '98}, May 14-16, Sinaia, 1998. p.62-67.

\bibitem[Dijkstra '65]{Dijkstra} E. W. Dijkstra (1965) Cooperating sequential processes. Technical Report EWD-123, Eindhoven University of Technology, Netherlands. Reprinted in {\em  Programming Languages} Academic Press, New York, 43--112. At: \url{https://www.cs.utexas.edu/users/EWD/transcriptions/EWD01xx/EWD123.html}


\bibitem[Dr\v ag\v anescu '84]{Draganescu84} Mihai Dr\v ag\v anescu: ``Information,
Heuristics, Creation", in Plauder, I. (ed): {\em Artificial Inteligence and Information Control System of Robots}, Elsevier Publishers B. V. (North-Holland), 1984.

%\bibitem[Dr\v{a}g\v{a}nescu '91]{} Mihai Dr\v{a}g\v{a}nescu, Gheorghe \c{S}tefan, Cornel Burileanu: {\em Electronica functional\v{a}}, Ed. Tehnic\v{a}, Bucure\c{s}ti, 1991 (in Roumanian).

%\bibitem[Einspruch '86]{} N. G. Einspruch ed.: {\em VLSI Electronics. Microstructure Science. vol. 14 : VLSI Design}, Academic Press, Inc., 1986.

%\bibitem[Einspruch '91]{} N. G. Einspruch, J. L. Hilbert: {\em Application Specific Integrated Circuits (ASIC) Technology}, Academic Press, Inc., 1991.

%\bibitem[Ercegovac '04]{Ercegovac} Milo\v s D. Ercegovac, Tom\' as Lang: {\em Digital Arithmetic}, Morgan Kaufman, 2004.

%\bibitem[Flynn '72]{Flynn} Flynn, M.J.: ``Some computer organization and their affectiveness", {\em IEEE Trans. Comp.} C21:9 (Sept. 1972), pp. 948-960.

\bibitem[Fortune '78]{Fortune} S. Fortune,  J. C. Wyllie (1978)  Parallelism in random access machines. {\em Conference Record of the Tenth Annual ACM Symposium on Theory of Computing},  114--118.

\bibitem[Goldschlage '82]{Goldschlage} L. M. Goldschlage (1982)  A universal interconnection pattern for parallel computers. {\em Journal of the ACM} {\bf 29}(4):1073--1086.

\bibitem[Gheolbanoiu '14]{Gheolbanoiu} Alexandru Gheolbanoiu, Dan Mocanu, Radu Hobincu, Lucian Petrica: ``Cellular Automaton pRNG with a Global Loop for Non-Uniform Rule Control", 18th International Conference on Ciruits, Systems, Communications and Computers (CSCC 2014), Santorini Island, Greece, July 17-21, 2014, vol. II, 415-420. Available at:  \url{http://www.europment.org/library/2014/santorini/bypaper/COMPUTERS/COMPUTERS2-15.pdf}

%\bibitem[Glushkov '66]{} V. M. Glushkov: {\em Introduction to Cybernetics}, Academic Press, 1966.

\bibitem[G\"{o}dels '31]{Godel} Kurt G\"{o}del: ``On Formally Decidable Propositions of
 Principia Mathematica and Related Systems I", reprinted in S. Fefermann et
 all.: {\em Collected Works I: Publications 1929 - 1936,} Oxford Univ. Press,
 New York, 1986.

%\bibitem[Hartley '95]{} Richard I. Hartley: {\em Digit-Serial Computation}, Kulwer Academic Pub., 1995.

%\bibitem[Hascsi '95]{} Zoltan Hascsi, Gheorghe \c{S}tefan: ``The Connex Content Addressable Memory ($C^{2}AM$)",  {\em Proceedings of the Twenty-first European Solid-State Circuits Conference}, Lille -France, 19-21 September 1995, pp. 422-425.

%\bibitem[Hascsi '96]{Hascsi} Zoltan Hascsi, Bogdan M\^{\i}\c tu, Mariana Petre, Gheorghe \c{S}tefan, ``High-Level Synthesis of an Enchanced Connex memory", in {\em Proceedings of the InternationalSemiconductor Conference}, Sinaia, October 1996, p. 163-166.

%\bibitem[Head '87]{} T. Head: ``Formal Language Theory and DNA: an Analysis of the Generative Capacity of Specific Recombinant Behaviours", in {\em Bull. Math. Biology}, 49, p. 737-759, 1987.

%\bibitem[Helbing 89']{} Walter A. Helbing, Veljko M. Milutinovic: ``Architecture and Design of a 32-bit GaAs Microprocessor", in [Milutinovic 89'].

\bibitem[Hennessy '07]{Hennessy} John L. Hennessy, David A. Patterson: {\em Computer Architecture: A Quantitative Approach}, Fourth Edition, Morgan Kaufmann, 2007.

\bibitem[Hennessy '19]{Hennessy1} J. L. Hennessy, D. A. Patterson (2019) {\it Computer Adchitecture. A Quantitative Approach. Sixth Edition}, Morgan Kaufmann.


%\bibitem[Hennie '68]{} F. C. Hennie: {\em Finite-State Models for Logical Machine}, John Wiley \& Sons, Inc., 1968.

\bibitem[Hill \& Reddi '21]{Hill} Mark D. Hill, Vijay Janapa Reddi: ``Accelerator-Level Parallelism", {\em Communications of the ACM}, {\bf 64}(12):36-38, 2021.

\bibitem[Hillis '85]{Hillis} W. D. Hillis: {\em The Connection Machine}, The MIT Press, Cambridge, Mass., 1985.

\bibitem[Hilbert 1900]{Hilbert} D. Hilbert, {\em Mathematical Problems. Lecture Delivered Before the International Congress of Mathematicians at Paris in 1900}.    Available at:  \url{https://www.ams.org/journals/bull/1902-08-10/S0002-9904-1902-00923-3/S0002-9904-1902-00923-3.pdf}

\bibitem[Hilbert \& Ackermann '28]{Hilbert28}    David Hilbert, Wilhelm Ackermann (1928). {\em Grundzüge der theoretischen Logik} (Principles of Mathematical Logic). Springer-Verlag.

%\bibitem[Kaeslin '01]{} Hubert Kaeslin: {\em Digital Integrated Circuit  Design}, Cambridge Univ. Press, 2008.

%\bibitem[Keeth '01]{Keeth} Brent Keeth, R. jacob Baker: {\em DRAM Circuit Design. A Tutorial}, IEEE Press, 2001.

\bibitem[Jouppi '17]{Jouppi}  Jouppi, N. P., Young, C., Patil, N., Patterson, D. {\em et al.} (2017) In-Datacenter Performance Analysis of a Tensor Processing Unit$^{TM}$, {\em 44th International Symposium on Computer Architecture (ISCA)}, Toronto, Canada, June 26.
    At: \url{https://drive.google.com/file/d/0Bx4hafXDDq2EMzRNcy1vSUxtcEk/view}

\bibitem[Kleene '36]{kleene} Stephen C. Kleene: ``General Recursive Functions of
 Natural Numbers", in {\em Math. Ann}., 112, 1936.

%\bibitem[Karim '08]{Karim} Mohammad A. Karim, Xinghao Chen: {\em Digital Design}, CRC Press, 2008.

\bibitem[Knuth '73]{Knuth} D. E. Knuth: {\em The Art of Programming. Sorting and Searching}, Addison-Wesley, 1973.

\bibitem[Kolmogorov '65]{Kolmogorov} A.A. Kolmogorov: ``Three Approaches to the
Definition of the Concept ``Quantity of Information" ``, in {\em
Probl. Peredachi Inform.}, vol. 1, pag. 3-11, 1965.

\bibitem[Kung '79]{Kung} H. T. Kung, C. E. Leiserson: ``Algorithms for VLSI processor arrays", in \cite{Mead}.

\bibitem[Ladner '80]{Ladner} R. E. Ladner, M. J. Fischer: ``Parallel prefix computation", {\em J. ACM}, Oct. 1980.

%\bibitem[Lindenmayer '68]{} Lindenmayer, A.: "Mathematical Models of Cellular Interactions in Development I, II", {\em Journal of Theor. Biology}, 18, 1968.

%\bibitem[Mali\c ta '06]{} Mihaela Mali\c ta, Gheorghe \c Stefan, Marius Stoian: ``Complex vs. Intensive in Parallel Computation", in {\em International Multi-Conference on Computing in the Global Information Technology - Challenges for the Next Generation of IT\&C - ICCGI}, 2006 Bucharest, Romania, August 1-3, 2006

%\bibitem[Mali\c ta '07]{} Mihaela Mali\c ta, Gheorghe \c Stefan, Dominique Thi\' ebaut: ``Not Multi-, but Many-Core: Designing Integral Parallel Architectures for Embedded Computation" in {\em International Workshop on Advanced Low Power Systems held in conjunction with 21st International Conference on Supercomputing} June 17, 2007 Seattle, WA, USA.

\bibitem[Mali\c ta '13]{malita13} Mihaela Mali\c ta, Gheorghe M. \c Stefan: ``Control Global Loops in Self-Organizing Systems", {\em ROMJIST}, Volume 16, Numbers 2–3, 2013, 177-191. \\ {\tt http://www.imt.ro/romjist/Volum16/Number16\_2/pdf/05-Malita-Stefan2.pdf}

%\bibitem[Markov '54]{} Markov, A. A.: "The Theory of Algorithms", {\em Trudy Matem. Instituta im V. A. Steklova}, vol. 42, 1954. (Translated from Russian by J. J. Schorr-kon, U. S. Dept. of Commerce, Office of Technical Services, no. OTS 60-51085, 1954)

\bibitem[Mead '79]{Mead} Carver Mead, Lynn Convay: {\em Introduction to VLSI Systems}, Addison-Wesley Pub, 1979.

%\bibitem[MicroBlaze]{MicroBlaze} *** {\em MicroBlaze Processor. Reference Guide.} posted at:\\ http://www.xilinx.com/support/documentation/sw\_manuals/xilinx14\_1/mb\_ref\_guide.pdf

\bibitem[Milutinovic 89']{Milutinovic} Veljko M. Milutinovic (ed.): {\em High-Level Language Computer Architecture}, Computer Science Press, 1989.

%\bibitem[Mindell '00]{mindell} Arnold Mindell: {\em Quantum Mind. The Edge Between Physics and Psychology}, Lao Tse Press, 2000.


%\bibitem[Minsky '67]{} M. L. Minsky: {\em Computation: Finite and Infinite Machine}, Prentice - Hall, Inc., 1967.

%\bibitem[M\^{\i}\c tu '00]{} Bogdan M\^{\i}\c tu, Gheorghe \c Stefan, ``Low-Power Oriented Microcontroller Architecture", in {\em CAS 2000 Proceedings}, Oct. 2000, Sinaia, Romania
\bibitem[Moore '11]{MooreC} Chuck Moore (2011) Data Processing in Exascale-class Computing Systems, {\em The Salishan Conference on
High Speed Computing} \url{https://www.lanl.gov/conferences/salishan/salishan2011/3moore.pdf}


\bibitem[Moore '65]{Moore} G. E. Moore (1965) Cramming more components onto integrated circuits, {\em Electronics}, {\bf 38}(8):114-117. %\url{https://www.alejandrobarros.com/wp-content/uploads/old/Articulo_original_G_Moore.pdf}

\bibitem[Moore '75]{Moore1} G. E. Moore (1975) Progress In Digital Integrated Electronics,  {\em IEEE, International Electron Devices Meeting. Technical Digest} pp.11-13.


\bibitem[Moto-Oka '82]{Moto} T. Moto-Oka (ed.): {\em Fifth Generation Computer Systems}, North-HollandPub. Comp., 1982.


\bibitem[Omondi '94]{Omondi} Amos R. Omondi: {\em Computer Arithmetic. Algorithm, Architecture and Implementation}, Prentice Hall, 1994.

%\bibitem[Palnitkar '96]{Palnitkar} Samir Palnitkar: {\em Verilog HDL. AGuide to Digital Design and Synthesis}, SunSoft Press, 1996.

\bibitem[Parberry 87]{Parberry} Ian Parberry: {\em Parallel Complexity Theory}. Research Notes in Theoretical Computer science. Pitman Publishing, London, 1987.

%\bibitem[Parberry 94]{} Ian Parberry: {\em Circuit Complexity and Neural Networks}, The MIT Presss, 1994.

\bibitem[Patterson '10]{Patterson} David A. Patterson (2010) The trouble with multicore. {\em IEEE Spectrum}, {\bf 47}(7):28-32.


\bibitem[Patterson '97]{Patterson1} David Patterson, et all, A case for intelligent RAM (1997) {\em  IEEE Micro} {\bf 17}(2)34:44.

\bibitem[Patterson '05]{Patterson2} David A. Patterson, John L.Hennessy: {\em Computer Organization \& Design. The Hardware / Software Interface}, Third Edition, Morgan Kaufmann, 2005.

\bibitem[Patterson '19]{Patterson3} David A. Patterson, John L.Hennessy: {\em Computer Organization \& Design. The Hardware / Software Interface}, Sixth Edition, Morgan Kaufmann, 2019.

%\bibitem[P\v aun '95a]{} P\v aun, G. (ed.): {\em Artificial Life. Grammatical Models}, Black Sea University Press, 1995.

%\bibitem[P\v aun '85]{Paun} A. P\v aun, Gh. \c Stefan, A. Birnbaum, V. Bistriceanu, ``DIALISP - experiment de structurare neconventionala a unei masini LISP", in {\em Calculatoarele electronice ale generatiei a cincea}, Ed. Academiei RSR, Bucuresti 1985. p. 160 - 165.

\bibitem[Post '36]{Post} Emil Post: ``Finite Combinatory Processes.
 Formulation I", in{\em The Journal of Symbolic Logic}, vol. 1, p. 103 -105, 1936.

 \bibitem[Pratt '74]{Pratt74}  V. R. Pratt, M. O. Rabin, L. J. Stockmeyer (1974)  A characterization of the power of vector machines. {\em Proceedings of STOC’1974}, pp. 122-134.

 \bibitem[Raina '16]{Raina}  G. Raina (2016) Deep Convolutional Netvork evaluation on the Xeon Phi: Where Subword Parallelism meets Many-Core, {\em Eindhoven University of Technology}, Master Thesis. At:  \url{https://repository.tue.nl/844256}

\bibitem[Redmon '16]{Redmon}      J. Redmon,  S. Divvala,  R. Girshick, A.   Farhadi (2016) You only look once: Unified, real-time object detection,  {\em Cornell Univ. Library}.

%\bibitem[Prince '99]{} Betty Prince: {\em High Performance Memories. New architecture DRAMs and SRAMs evolution anad function}, John Wiley \& Sons, 1999.

%\bibitem[Rafiquzzaman '05]{Rafiquzzaman} Mohamed Rafiquzzaman: {\em Fundamentals of Digital Logic and Microcomputer Design}, Fifth Edition,  Wiley -- Interscience, 2005.

%\bibitem[Salomaa '73]{} Arto Salomaa: {\em Formal Languages}, Academic Press, Inc., 1973.

%\bibitem[Salomaa '81]{} Arto Salomaa: {\em Jewels of Formal Language Theory}, Computer Science Press, Inc., 1981.

\bibitem[Savage '87]{Savage} John Savage: {\em The Complexity of Computing}, Robert E. Krieger Pub. Comp., 1987.

%\bibitem[Shankar '89]{} R. Shankar, E. B. Fernandez: {\em VLSI Computer Architecture}, Academic Press, Inc., 1989.

\bibitem[Shannon '38]{Shannon '38} C. E. Shannon: ``A Symbolic Analysis of Relay and Switching Circuits", {\em Trans. American Institute of Electrical Engineers}, {\bf  57}(12):713-723, 1938. (The paper is an abstract of the thesis presented in 1937 at MIT for the degree of master in science.)

\bibitem[Shannon '48]{Shannon48} C. E. Shannon: ``A Mathematical Theory of Communication", {\em Bell System Tech. J.}, Vol. 27, 1948.

\bibitem[Shannon '56]{Shannon56} C. E. Shannon: ``A Universal Turing Machine with Two Internal States", in {\em Annals of Mathematics Studies, No. 34: Automata Studies}, Princeton Univ. Press, pp 157-165, 1956.

%\bibitem[Sharma '97]{Sharma} Ashok K. Sharma: {\em Semiconductor Memories. Techology, Testing, and Reliability}, Wiley -- Interscience, 1997.

%\bibitem[Sharma '03]{} Ashok K. Sharma: {\em Advanced Smiconductor Memories. Architectures, Designs, and Applications}, Whiley-Interscience, 2003.

\bibitem[Solomonoff '64]{Solomonoff} R. J. Solomonoff: ``A Formal Theory of Inductive Inference", in {\em Information and Control}, vol. 7, no. 2 , pag. 224-254, 1964.

\bibitem[Spira '71]{Spira} P. M. Spira: ``On time-Hardware Complexity Tradeoff for Boolean Functions", in {\em Preceedings of Fourth Hawaii International Symposium on System Sciences}, pp. 525-527, 1971.

%\bibitem[Stoian '07]{} Marius Stoian, Gheorghe \c Stefan: ``Stacks orFile-Registers in Cellular Computing?", in {\em CAS, Sinaia 2007}.

\bibitem[Streinu '85]{Streinu} Ileana Streinu: ``Explicit Computation of an Independent G\"{o}del Sentence", in {\em Recursive Functions Theory Newsletter}, June 1985.

\bibitem[Stefan '20]{Stefan1}    G. M. \c Stefan (2020) {\em Composition is the only independent rule in Kleene's model of partial recursive functions}, At: \url{http://users.dcae.pub.ro/~gstefan/2ndLevel/composition.html}

%\bibitem[\c Stefan '97]{DStefan} Denisa \c Stefan, Gheorghe \c Stefan, ``Bi-thread Microcontroller as Digital Signal Processor", in {\em CAS '97 Proceedings, 1997 International Semiconductor Conference}, October 7 -11, 1997, Sinaia, Romania.

%\bibitem[\c Stefan '99]{DStefan} Denisa \c Stefan, Gheorghe \c Stefan: ``A Procesor Network without Interconnectio Path", in {\em CAS 99 Proceedings, Oct., 1999}, Sinaia, Romania. p. 305-308.

%\bibitem[\c Stefan '80]{}  Gheorghe \c Stefan: {\em LSI Circuits for Processors}, Ph.D. Thesis (in Roumanian), Politechnical Institute of Bucharest, 1980.

%\bibitem[\c Stefan '83]{Stefan83}  Gheorghe \c Stefan: ``Structurari neechilibrate in sisteme de prelucrare a informatiei", in {\em Inteligenta artificiala si robotica}, Ed. Academiei RSR, Bucuresti, 1983. p. 129 - 140.

%\bibitem[\c Stefan '83]{}  Gheorghe \c Stefan, et al.: {\em Circuite integrate digitale}, Ed. Did. si Ped., Bucuresti, 1983.

%\bibitem[\c Stefan '84]{}   Gheorghe \c Stefan,  et al.: ``DIALISP - a LISP Machine", in {\em Proceedings of the ACM Symposium on LISP and Functional Programming}, Austin, Texas, Aug. 1984. p. 123 - 128.

%\bibitem[\c Stefan '85]{Stefan85}  Gheorghe \c Stefan, A. P\v aun, ``Compatibilitatea functie - structura ca mecanism al evolutiei arhitecturale", in {\em Calculatoarele electronice ale generatiei a cincea}, Ed. Academiei RSR, Bucuresti, 1985. p. 113 - 135.

%\bibitem[\c Stefan '85a]{Stefan85a}  Gheorghe \c Stefan, V. Bistriceanu, A. P\v aun, ``Catre un mod natural de implementare a LISP-ului", in {\em Sisteme cu inteligenta artificiala}, Ed. Academiei Romane, Bucuresti, 1991 (paper at {\em Al doilea simpozion national deinteligenta artificiala}, Sept. 1985). p. 218 224.

%\bibitem[\c Stefan '86]{Stefan86}  Gheorghe Stefan, M. Bodea, ``Note de lectura la volumul lui T. Blakeslee: Proiectarea cu circuite MSI si LSI", in {\em T. Blakeslee: Prioectarea cu circuite integrate MSI si LSI}, Ed. Tehnica, Bucuresti, 1986 (translated from English by M. Bodea, M. Hancu, Gh. Stefan). p. 338 - 364.

%\bibitem[\c Stefan '86a]{Stefan86a}  Gheorghe Stefan, ``Memorie conexa" in {\em CNETAC 1986} Vol. 2, IPB, Bucuresti, 1986, p. 79 - 81.

\bibitem[\c Stefan '91]{Stefan91}  Gheorghe \c Stefan: {\em Funcție și structură în sistemele digitale} (Function and Structure in Digital Systems), Ed. Academiei Române, 1991.

%\bibitem[\c Stefan '91]{}   Gheorghe \c Stefan,  Dr\v aghici, F.: ``Memory Management Unit - a New Principle for LRU Implementation", {\em Proceedings of 6th Mediterranean Electrotechnical Conference}, Ljubljana, Yugoslavia, May 1991, pp. 281-284.

%\bibitem[\c Stefan '93]{}  Gheorghe \c Stefan: {\em Circuite integrate digitale}. Ed. Denix, 1993.

%\bibitem[\c Stefan '95]{} Gheorghe \c Stefan,   Mali\c ta, M.: ``The Eco-Chip: A Physical Support for Artificial Life Systems",  {\em Artificial Life. Grammatical Models}, ed. by Gh. P\v aun, Black Sea University Press, Bucharest, 1995, pp. 260-275.

%\bibitem[\c Stefan '96]{}  Gheorghe \c Stefan, Mihaela Mali\c ta: ``Chaitin's Toy-Lisp  on Connex Memory Machine",  {\em Journal of Universal Computer Science}, vol. 2, no. 5, 1996, pp. 410-426.

%\bibitem[\c Stefan '97]{Stefan97}  Gheorghe \c Stefan, Mihaela Mali\c ta: ``DNA Computing with the Connex Memory", in {\em RECOMB 97 First International Conference on Computational Molecular Biology}. January 20 - 23, Santa Fe, New Mexico, 1997. p. 97-98.

%\bibitem[\c{S}tefan '97a]{Stefan97a}  Gheorghe \c Stefan, Mihaela Mali\c ta: `` The Splicing Mechanism and the Connex Memory", {\em Proceedings of the 1997 IEEE International Conference on Evolutionary Computation}, Indianapolis, April 13 - 16, 1997. p. 225-229.
 %\bibitem[\c Stefan '98]{Stefan}  Gheorghe \c Stefan, ``Silicon or Molecules? What's the Best for Splicing", in Gheorghe P\v aun (ed.): {\em Computing with Bio-Molecules. Theory and Experiments.} Springer, 1998. p. 158-181

\bibitem[\c Stefan '98a]{Stefana}  Gheorghe \c Stefan,  `` ``Looking for the Lost Noise" ", in {\em CAS '98 Proceedings, Oct. 6 - 10, 1998}, Sinaia, Romania. p.579 - 582. \\ {\tt http://arh.pub.ro/gstefan/CAS98.pdf}

%\bibitem[\c Stefan '98b]{Stefanb}  Gheorghe \c Stefan, ``The Connex Memory: A Physical Support for Tree / List Processing" in {\em The Roumanian Journal of Information Science and Technology}, Vol.1, Number 1, 1998, p. 85 - 104.

%\bibitem[\c Stefan '98]{}  Gheorghe \c Stefan, Robrt Benea: ``Connex Memories \& Rewrieting Systems", in {\em MELECON '98}, Tel-Aviv, May 18 -20, 1998.

%\bibitem[\c Stefan '99]{}  Gheorghe \c  Stefan, Robert Benea: ``Experimente in info cu acizi nucleici", in M. Dr\v ag\v anescu, \c Stefan Tr\v ausan-Matu (eds): {\em Natura realitatii fizice si a informatiei}, Editura Tempus, 1999.

%\bibitem[\c Stefan '99a]{}  Gheorghe \c Stefan: ``A Multi-Thread Approach in Order to Avoid Pipeline Penalties", in {\em Proceedings of 12th International Conference on Control Systems and Computer Science}, Vol. II, May 26-29, 1999, Bucharest, Romania. p. 157-162.

%\bibitem[\c Stefan '00]{}  Gheorghe \c Stefan: ``Parallel Architecturing starting from NaturalComputational Models", in {\em Proceedings of the Romanian Academy}, Series A: Mathematics, Physics, Technical Sciences, Information Science, vol. 1, no. 3 Sept-Dec 2000.

%\bibitem[\c Stefan '01]{}  Gheorghe \c Stefan, Dominique Thi\' ebaut, ``Hardware-Assisted String-Matching Algorithms", in {\em WABI 2001, 1st Workshop on Algorithms in BioInformatics, BRICS}, University of Aarhaus, Danemark, August 28-31, 2001.

\bibitem[\c Stefan '04]{Stefan04}  Gheorghe \c Stefan, Mihaela Mali\c ta: ``Granularity and Complexity in Parallel Systems", in {\em Proceedings of the 15 IASTED International Conf, 2004}, Marina Del Rey, CA, ISBN 0-88986-391-1, pp.442-447.

%\bibitem[\c Stefan '06]{}  Gheorghe \c Stefan: ``Integral Parallel Computation", in {\em Proceedings of the Romanian Academy, Series A: Mathematics, Physics, Technical Sciences, Information Science}, vol. 7, no. 3 Sept-Dec 2006, p.233-240.

\bibitem[\c Stefan '06a]{Stefan6a}  Gheorghe \c Stefan: ``A Universal Turing Machine with Zero Internal
States", in {\em Romanian Journal of Information Science and
Technology}, {\bf 9}(3):227-243,  2006.

%\bibitem[\c Stefan '06b]{}  Gheorghe \c Stefan: ``The CA1024: SoC with Integral Parallel Architecture for HDTV Processing", invited paper at {\em 4th International System-on-Chip (SoC) Conference \& Exhibit}, November 1 \& 2, 2006, Radisson Hotel Newport Beach, CA

%\bibitem[\c Stefan '06c]{}  Gheorghe \c Stefan, Anand Sheel, Bogdan M\^{\i}\c tu, Tom Thomson, Dan Tomescu: ``The CA1024: A Fully Programmable System-On-Chip for Cost-Effective HDTV Media Processing", in {\em Hot Chips: A Symposium on High Performance Chips}, Memorial Auditorium, Stanford University, August 20 to 22, 2006.

%\bibitem[\c Stefan '06d]{}  Gheorghe \c Stefan: ``The CA1024: A Massively Parallel Processor for Cost-Effective HDTV", in {\em SPRING PROCESSOR FORUM: Power-Efficient Design}, May 15-17, 2006, Doubletree Hotel, San Jose, CA.

%\bibitem[\c Stefan '06e]{}  Gheorghe \c Stefan: ``The CA1024: A Massively Parallel Processor for Cost-Effective HDTV", in {\em SPRING PROCESSOR FORUM JAPAN}, June 8-9, 2006, Tokyo.

%\bibitem[\c Stefan '07]{}  Gheorghe \c Stefan: ``Membrane Computing in Connex Environment", invited paper at {\em 8th Workshop on Membrane Computing (WMC8)} June 25-28, 2007 Thessaloniki, Greece

%\bibitem[\c Stefan '07a]{}  Gheorghe \c Stefan, Marius Stoian: ``The efficiency of the register file based architectures in OOP languages era", in {\em SINTES13} Craiova, 2007.

%\bibitem[\c Stefan '07b]{}  Gheorghe \c Stefan: ``Chomsky's Hierarchy \& a Loop-Based Taxonomy for Digital Systems", in {\em Romanian Journal of Information Science and Technology} vol. 10, no. 2, 2007.

\bibitem[\c Stefan '14]{Stefan14} Gheorghe M. Stefan, Mihaela Malita: ``Can One-Chip Parallel Computing Be Liberated From Ad Hoc Solutions? A Computation Model Based Approach and Its Implementation", {\em 18th International Conference on Circuits, Systems, Communications and Computers (CSCC 2014)}, Santorini Island, Greece, July 17-21, 2014, 582-597. \\ {\tt  http://www.inase.org/library/2014/santorini/bypaper/COMPUTERS/COMPUTERS2-42.pdf}

\bibitem[\c Stefan '21a]{Stefan21a} Gheorghe M. Stefan:  ``Let's consider Moore's law in its entirety", CAS 2021 PROCEEDINGS : 2021 International Semiconductor Conference : 44nd edition, , pp. 3-10, October 6-8, Sinaia, Romania.

\bibitem[\c Stefan '21]{Stefan21} Gheorghe M. Stefan:  ``Pseudo-Reconfigurable Systems", ROMJIST, {\bf 24}(12):366-383, 2021.

%\bibitem[Sutherland '02]{Sutherland} Stuart Sutherland: {\em Verilog 2001. A Guide to the New Features of the Verilog Hardware Description Language}, Kluwer Academic Publishers, 2002.

%\bibitem[Tabak '91]{} D. Tabak: {\em Advanced Microprocessors}, McGrow- Hill, Inc., 1991.

%\bibitem[Tanenbaum '90]{} A. S. Tanenbaum: {\em Structured Computer Organisation} third edition, Prentice-Hall, 1990.

%\bibitem[Thi\' ebaut '06]{} Dominique Thi\' ebaut, Gheorghe \c Stefan, Mihaela Mali\c ta: ``DNA search and the Connex technology" in {\em International Multi-Conference on Computing in the Global Information Technology - Challenges for the Next Generation of IT\&C - ICCGI}, 2006 Bucharest, Romania, August 1-3, 2006

%\bibitem[Tokheim '94]{Tokheim} Roger L. Tokheim: {\em Digital Principles}, Third Edition, McGraw-Hill, 1994.

\bibitem[Turing '36]{Turing} Alan M. Turing: ``On computable Numbers with an
 Application to the Eintscheidungsproblem", in {\em Proc. London Mathematical
 Society,} 42 (1936), 43 (1937).

%\bibitem[Vahid '06]{Vahid} Frank Vahid: {\em Digital Design}, Wiley, 2006.

\bibitem[von Neumann '45]{von Neumann} John von Neumann: ``First Draft of a Report on
 the EDVAC", reprinted in {\em IEEE Annals of the History of Computing,}
 Vol. 5, No. 4, 1993.

%\bibitem[Uyemura '02]{Uyemura} John P. Uyemura: {\em CMOS Logic Circuit Design}, Kluver Academic Publishers, 2002.

%\bibitem[Ward '90]{} S. A. Ward, R. H. Halstead: {\em Computation Structures}, The MIT Press, McGraw-Hill Book Company, 1990.

%\bibitem[Wedig '89]{} Robert G. Wedig: ``Direct Correspondence Architectures: Principles, Architecture, and Design" in [Milutinovic '89].

\bibitem[Waksman '68]{Waksman} Abraham Waksman, "A permutation network," in {\em J. Ass. Comput. Mach.}, vol. 15, pp. 159-163, Jan. 1968.

%\bibitem[webRef\_1]{webRef1} \url{http://www.fpga-faq.com/FAQ_Pages/0017_Tell_me_about_metastables.htm}

%\bibitem[webRef\_2]{webRef2} \url{http://www.fpga-faq.com/Images/meta_pic_1.jpg}


\bibitem[webRef\_3]{webRef3} \url{http://www.aoki.ecei.tohoku.ac.jp/arith/mg/algorithm.html#fsa_pfx}


%\bibitem[Weste '94]{Weste} Neil H. E. Weste, Kamran Eshraghian: {\em Principle of CMOS VLSI Design. ASystem Perspective}, Second Edition, Addisson Wesley, 1994.

%\bibitem[Wolfram '02]{Wolfram} Stephen Wolfram: {\em A New Kind of Science}, Wolfram Media, Inc., 2002.

\bibitem[Zurada '95]{Zurada} Jacek M. Zurada: {\em Introductin to Artificial Neural network}, PWS Pub. Company, 1995.

%\bibitem[Yanushkevich '08] {Yanushkevich} Svetlana N. Yanushkevich, Vlad P. Shmerko: {\em Introduction to Logic Design}, Press, 2008.



\end{thebibliography}
